\chapter{Conclusiones y trabajo futuro}
\label{cap:conclusiones}

\section{Conclusiones}
\textit{(Redactar a partir de los resultados obtenidos.)} En este Trabajo de Fin
de Grado se ha estudiado la aplicación del Análisis Topológico de Datos a la
clasificación y visualización de mamografías digitales del conjunto
\mbox{CBIS-DDSM}. Las principales conclusiones, en relación con los objetivos
planteados, son:

\begin{itemize}
  \item Se ha revisado el estado del arte en análisis de mamografías y en TDA
        aplicado a imagen médica (O1).
  \item Se ha construido un \textit{pipeline} reproducible de preparación de
        datos con partición por paciente que evita la fuga de datos (O2).
  \item Se han extraído y evaluado descriptores topológicos como características
        para la clasificación benigno/maligno (O3, O4).
  \item Se ha comparado el enfoque topológico con un modelo de \textit{deep
        learning} de referencia y se ha explorado su combinación (O5).
  \item Se han desarrollado visualizaciones que mejoran la interpretabilidad de
        los resultados (O6).
\end{itemize}

\section{Grado de cumplimiento de los objetivos}
Se valora el grado de cumplimiento de cada objetivo específico y las
desviaciones respecto a la planificación inicial.

\section{Limitaciones}
Se discuten las limitaciones del trabajo: tamaño y sesgos del conjunto de datos,
coste computacional de la homología persistente, y alcance de la validación
clínica.

\section{Trabajo futuro}
Como líneas de trabajo futuro se proponen:
\begin{itemize}
  \item Extender el estudio a otras filtraciones y descriptores topológicos
        (\textit{multiparameter persistence}).
  \item Integrar el TDA directamente en arquitecturas de \textit{deep learning}
        (capas topológicas diferenciables).
  \item Validar el enfoque en otros conjuntos de datos (INbreast, datos propios).
  \item Realizar una evaluación con radiólogos para medir la utilidad clínica de
        las visualizaciones.
\end{itemize}
